\section{Committing Security and Context Discoverability}
\label{sec:cmtds3}

This section proves the random-oracle CMTDs-3 bound that the headline
CMTDs-4 statement (\Cref{thm:cmtds4}) consumes via the trivial
promotion of \Cref{app:committing-transfer}. The public-permutation
CMTDs-3 statement follows by the flat-sponge lift of
\Cref{app:flat-sponge-lift}. The same root-tag argument, with the
$s$-way collision replaced by a preimage to a fixed target tag, gives
the random-oracle context-discoverability bound
(\Cref{thm:cdy-ro}) underlying the headline CDY statement
(\Cref{thm:cdy}). A revealed-target variant of the same preimage
argument gives the random-oracle second-preimage bound
(\Cref{cor:spr-ro}) underlying the headline second-preimage corollary
(\Cref{cor:spr}), the realized-tag companion to context discoverability.

\begin{theorem}[Random-Oracle CMTDs-3]
\label{thm:cmtds3-ro}
For every $s \geq 2$ and every $s$-way CMTDs-3
adversary $\mathcal{B}$ in the \TW{} random-oracle model
(\Cref{sec:ro-model}),
\[
  \Advcmtds{3}_\RO(\mathcal{B})
  \;\leq\; \RootColl_{\tauroot}(\qRO(\mathcal{B}), s)
  = \frac{\binom{\qRO(\mathcal{B}) + s}{s}}{2^{\tauroot (s - 1)}}.
\]
\end{theorem}

\begin{proof}
Unlike the privacy and INT-CTXT proofs, the CMTDs-3 analysis has no
$\badkey$ event: all $s$ keys are
adversary-supplied, so a direct query to a keyed transcript is one of
$\mathcal{B}$'s visible $\ROflat$ queries and is counted in
$\qRO(\mathcal{B})$.

Run the random-oracle CMTDs-3 game (\Cref{sec:ro-model}) and define a
candidate sequence for every execution. During $\mathcal{B}$'s direct
oracle phase, when $\mathcal{B}$ makes a direct query to an input $Y$
that is not already in the sequence, append $Y$ immediately before
answering the query if $\KeyedTWOut(Y)$ decodes to a root tag
output-point class, namely $\Rmsg^+$ (the single-chunk root tag, with
the aggregation phase elided when $n = 0$) or $\Rtag$ (the multi-chunk
root tag). This predicate depends on $Y$ alone through the exact decoder
of \Cref{lem:bridge-decode}, which ignores the requested output length,
so both the decision to append and the string $Y$ are determined before
the query's $\ROflat$ lookup. By Output-Point Separation
(\Cref{cor:output-sep}) each input decodes to at most one output-point
class, so each distinct direct query contributes at most one candidate.
Since there is no construction oracle before $\mathcal{B}$ outputs its
tuples, each direct-query candidate is first touched at its append step.

After $\mathcal{B}$ outputs $C \concat T$ and tuples
$(K_i,N_i,A_i,M_i)_{i=1}^s$, stop appending challenger candidates if
the syntax, message-length, or pairwise-distinct-triple check fails.
Otherwise, as the
challenger evaluates the $s$ deterministic decryption checks, let
$R_i$ be the final root transcript reached by
$\Dec_{K_i}^\RO(N_i,A_i,C \concat T)$ in check $i$, and append
$\flatmap(R_i)$ immediately before the root tag oracle answer is
sampled, unless that input is already in the sequence. If a challenger
input was already touched by a direct query, then
\Cref{lem:bridge-decode} and Output-Point Separation
(\Cref{cor:output-sep}) imply that the earlier query was a final-root
candidate and the input is already in the sequence. Otherwise,
\Cref{cor:output-sep} excludes all earlier challenger lookups outside
the final-root class, and Opening Triple Separation
(\Cref{cor:triple-sep}) excludes earlier final-root lookups from the
other checked triples. Thus no earlier lookup of any kind has touched
a newly appended candidate's input. The sequence therefore satisfies
the predictability and unrevealed-value premises of Adaptive Candidate
Collision (\Cref{lem:adaptive-candidate-collision}), and its length is
at most $\qRO(\mathcal{B}) + s$.

On a winning execution, the challenger has accepted the syntax,
message-length, and pairwise-distinct-triple checks, so the tuples are
valid and the $(K_i,N_i,A_i)$ triples are pairwise distinct. By
\Cref{cor:triple-sep}, the bridged final-root inputs
\[
  \flatmap(R_1),\dots,\flatmap(R_s)
\]
are pairwise distinct. Thus leaf tag collisions cannot merge the root
inputs relevant to a CMTDs-3 win. Since every decryption check
accepts, the
$\tauroot$-bit projections at the $s$ distinct challenger final-root
inputs all equal the supplied tag $T$:
\[
  \bigl\lfloor\ROflat(\flatmap(R_1))\bigr\rfloor_{\tauroot}
  =
  \cdots
  =
  \bigl\lfloor\ROflat(\flatmap(R_s))\bigr\rfloor_{\tauroot}
  =
  T .
\]
Hence the win event is contained in the $s$-way candidate-collision
event. Applying \Cref{lem:adaptive-candidate-collision} with
$Q=\qRO(\mathcal{B})$ and $\tau=\tauroot$ gives
\[
  \Advcmtds{3}_\RO(\mathcal{B})
  \leq
  \binom{\qRO(\mathcal{B}) + s}{s} \cdot 2^{-\tauroot(s-1)}
  =
  \RootColl_{\tauroot}(\qRO(\mathcal{B}), s).
\]
\end{proof}

\begin{theorem}[Random-Oracle CDY]
\label{thm:cdy-ro}
For every query-independent target tag $T^* \in \bits^{\tauroot}$ and
every context-discoverability adversary $\mathcal{B}$ in the \TW{}
random-oracle model (\Cref{sec:ro-model}),
\[
  \Advcdy_\RO(\mathcal{B})
  \;\leq\; \RootPre_{\tauroot}(\qRO(\mathcal{B}))
  = \frac{\qRO(\mathcal{B}) + 1}{2^{\tauroot}}.
\]
\end{theorem}

\begin{proof}
As in \Cref{thm:cmtds3-ro}, the CDY analysis has no $\badkey$ event:
the single key is adversary-supplied, so a direct query to a keyed
transcript is one of $\mathcal{B}$'s visible $\ROflat$ queries and is
counted in $\qRO(\mathcal{B})$. Fix the query-independent target
$T^*$, which is therefore independent of $\ROflat$.

Run the random-oracle CDY game (\Cref{sec:ro-model}) and build the
root-tag candidate sequence exactly as in the direct-oracle phase of
\Cref{thm:cmtds3-ro}: each distinct direct query whose input decodes to
a root tag output-point class ($\Rmsg^+$ or $\Rtag$) is appended before
its $\ROflat$ lookup and first touched at its append step.

After $\mathcal{B}$ outputs the ciphertext body $C$ and the single
context $(K,N,A)$, stop appending if the syntax or ciphertext-validity
precheck on $C \concat T^*$ fails. Otherwise let $R$ be the final root
transcript reached by the one deterministic decryption check
$\Dec_{K}^\RO(N,A,C \concat T^*)$, and append $\flatmap(R)$ immediately
before its root tag oracle answer is sampled, unless that input is
already in the sequence. If $\flatmap(R)$ was already touched by a
direct query, then \Cref{lem:bridge-decode} and Output-Point Separation
(\Cref{cor:output-sep}) imply the earlier query was a final-root
candidate and the input is already in the sequence; otherwise
\Cref{cor:output-sep} excludes all earlier challenger lookups outside
the final-root class. Unlike \Cref{thm:cmtds3-ro}, there is a single
context, so no appeal to Opening Triple Separation
(\Cref{cor:triple-sep}) is needed: no pairwise distinctness among
final-root inputs is required. Thus no earlier lookup of any kind has
touched the newly appended candidate's input. The sequence therefore
satisfies the predictability and unrevealed-value premises of Adaptive
Candidate Preimage (\Cref{lem:adaptive-candidate-preimage}), and its
length is at most $\qRO(\mathcal{B}) + 1$.

On a winning execution the precheck passes and the single decryption
check accepts. By \Cref{alg:dec}, acceptance is the constant-time
equality $T' = T^*$, where $T'$ is the $\tauroot$-bit projection at the
final root transcript $R$; in the random-oracle model
$T' = \bigl\lfloor \ROflat(\flatmap(R)) \bigr\rfloor_{\tauroot}$. Hence
the win event is contained in the event that some candidate's
$\tauroot$-bit projection equals the fixed target $T^*$. Applying
\Cref{lem:adaptive-candidate-preimage} case (i) with $Q = \qRO(\mathcal{B})$,
$\tau = \tauroot$, and $t = T^*$ gives
\[
  \Advcdy_\RO(\mathcal{B})
  \leq
  (\qRO(\mathcal{B}) + 1) \cdot 2^{-\tauroot}
  =
  \RootPre_{\tauroot}(\qRO(\mathcal{B})).
\]
\end{proof}

\begin{corollary}[Random-Oracle Second-Preimage]
\label{cor:spr-ro}
For every query-independent honest opening $\omega_0 = (K_0, N_0, A_0,
M_0)$ in the \TW{} encryption domain and every second-preimage adversary
$\mathcal{B}$ in the \TW{} random-oracle model (\Cref{sec:ro-model}),
\[
  \Advspr_\RO(\mathcal{B})
  \;\leq\; \RootPre_{\tauroot}(\qRO(\mathcal{B}))
    + \frac{\nleaf(\mathcal{B})}{2^{\tauleaf}}
  = \frac{\qRO(\mathcal{B}) + 1}{2^{\tauroot}}
    + \frac{\nleaf(\mathcal{B})}{2^{\tauleaf}}.
\]
\end{corollary}

\begin{proof}
As in \Cref{thm:cdy-ro}, the analysis has no $\badkey$ event: both the
honest opening's key $K_0$ and the forged context's key $K_1$ are
adversary-supplied, so a direct query to a keyed transcript is one of
$\mathcal{B}$'s visible $\ROflat$ queries and is counted in
$\qRO(\mathcal{B})$. Fix $\omega_0$.

Run the random-oracle second-preimage game (\Cref{sec:ro-model}). Before
$\mathcal{B}$ is run, the challenger computes the honest target
$C_0 \concat T = \Enc_{K_0}^\RO(N_0, A_0, M_0)$; its final root tag
lookup is at the address $X_0 = \flatmap(R_0)$, where $R_0$ is the final
root transcript of the honest encryption, and it fixes the realized
target $T = \lfloor \ROflat(X_0) \rfloor_{\tauroot}$. By Output-Point
Separation (\Cref{cor:output-sep}), the only honest-encryption lookup in
a root tag output-point class is the one at $X_0$: its remaining lookups
(root and leaf initialization, associated-data, non-final and leaf
message, and aggregation points) lie in other classes and are distinct
from every root tag address.

Build the root-tag candidate sequence as in the direct-oracle phase of
\Cref{thm:cmtds3-ro}, with the single change that the target address
$X_0$ is excluded: a direct query to $Y \neq X_0$ whose input decodes to
a root tag output-point class is appended before its $\ROflat$ lookup,
while a direct query to $X_0$ itself only re-reads the already-revealed
target and is not appended.

After $\mathcal{B}$ outputs the ciphertext body $C_1$ and the single
context $(K_1, N_1, A_1)$, stop appending if the syntax, ciphertext
validity, or distinctness precheck on $(C_1, K_1, N_1, A_1)$ fails.
Otherwise let $R_1$ be the final root transcript reached by the one
deterministic decryption check $\Dec_{K_1}^\RO(N_1, A_1, C_1 \concat T)$,
and append $\flatmap(R_1)$ immediately before its root tag oracle answer
is sampled, unless that input is already in the sequence. On a winning
execution the distinctness check passes, so
$(C_1, K_1, N_1, A_1) \neq (C_0, K_0, N_0, A_0)$.

Let $\badleaf$ be the event that the decryption check generates a final
leaf transcript distinct from the honest leaf transcript at the same leaf
index but with an equal $\tauleaf$-bit leaf tag. The honest leaf tags of
$\omega_0$ are fixed before $\mathcal{B}$ runs, so they form a revealed
target set; the candidate leaf transcripts are the decryption check's
leaf transcripts at honest-occupied leaf indices that are distinct from
the honest leaf transcript sharing their leaf index, at most
$\nleaf(\mathcal{B})$ in all. A check leaf at a leaf index with no honest
leaf transcript can never satisfy $\badleaf$---a same-index collision with
an honest leaf---so restricting to honest-occupied indices loses nothing.
Each candidate check leaf is construction-generated, hence predictable,
and differs from the honest leaf transcript at its occupied index, so by
\Cref{lem:transcript-inj} its bridged address is distinct from that
index's revealed target address; the candidates therefore satisfy the
predictability and address-separation premises of Revealed-Target Leaf
Preimage, and the candidate set is precisely the $\badleaf$ event. By
Revealed-Target Leaf Preimage
(\Cref{cor:revealed-target-leaf-preimage}),
\[
  \Pr[\badleaf]
  \;\leq\;
  \frac{\nleaf(\mathcal{B})}{2^{\tauleaf}}.
\]
On executions where $\badleaf$ does not occur, a distinct opening forces
$R_1 \neq R_0$: if the triples differ,
$(K_1, N_1, A_1) \neq (K_0, N_0, A_0)$, Opening Triple Separation
(\Cref{cor:triple-sep}) gives $\flatmap(R_1) \neq \flatmap(R_0)$
directly; if instead the triples agree and $C_1 \neq C_0$, suppose
$R_1 = R_0$. Equality of the flattened final-root inputs recovers the same
positional aggregation string
$T_1 \concat \cdots \concat T_n \concat \langle n \rangle_8$
(\Cref{lem:transcript-inj}), so the check and honest computations share a
leaf tag at every leaf index. Absence of $\badleaf$ then forces the
index-aligned check and honest leaf transcripts to coincide, so
\Cref{lem:transcript-inj} recovers the same absorbed ciphertext
$C_1 = C_0$, a contradiction; hence $R_1 \neq R_0$. Recall
$X_0 = \flatmap(R_0)$. In the first sub-case Opening Triple Separation
already gives the address inequality
$\flatmap(R_1) \neq \flatmap(R_0) = X_0$ directly; in the second,
$R_1 \neq R_0$, so injectivity of $\flatmap$ (\Cref{lem:bridge-decode})
gives $\flatmap(R_1) \neq \flatmap(R_0) = X_0$. The winning candidate
therefore lies at an address distinct from the target address $X_0$. As in \Cref{thm:cdy-ro}, if $\flatmap(R_1)$ was already touched by
a direct query then \Cref{cor:output-sep} makes that earlier query a
final-root candidate already in the sequence; otherwise \Cref{cor:output-sep}
excludes all earlier lookups outside the final-root class, and the honest
encryption's only root tag lookup is at the excluded address $X_0$. Thus
no earlier lookup has revealed the newly appended candidate's value. The
sequence therefore satisfies the predictability and unrevealed-value
premises of Adaptive Candidate Preimage
(\Cref{lem:adaptive-candidate-preimage}), every candidate is distinct from
$X_0$, and its length is at most $\qRO(\mathcal{B}) + 1$.

On a winning execution the single decryption check accepts, so by
\Cref{alg:dec} the $\tauroot$-bit projection at the final root transcript
$R_1$ equals the realized target:
\[
  \bigl\lfloor \ROflat(\flatmap(R_1)) \bigr\rfloor_{\tauroot}
  = T
  = \bigl\lfloor \ROflat(X_0) \bigr\rfloor_{\tauroot}.
\]
Hence, on executions where $\badleaf$ does not occur, the win event is
contained in the event that some candidate's $\tauroot$-bit projection
equals the revealed target $T$. Applying \Cref{lem:adaptive-candidate-preimage}
case (ii) with $Q = \qRO(\mathcal{B})$, $\tau = \tauroot$, distinguished
address $X_0$, and revealed target $T$ bounds that event by
$\RootPre_{\tauroot}(\qRO(\mathcal{B}))$. Combining with the $\badleaf$
bound above,
\[
  \Advspr_\RO(\mathcal{B})
  \leq
  \RootPre_{\tauroot}(\qRO(\mathcal{B}))
  + \frac{\nleaf(\mathcal{B})}{2^{\tauleaf}}.
\]
Since the bound is uniform in $\omega_0$, it holds for the maximum
defining $\Advspr_\RO$.
\end{proof}
